\documentclass[12pt]{article}
\usepackage{graphicx}
\usepackage{float}

\usepackage[T2A]{fontenc}
\usepackage[utf8]{inputenc}
\usepackage[russian]{babel}

\usepackage{amsmath, amsfonts, amssymb}
\usepackage{amsthm}
% \usepackage{mathpazo}
% \usepackage{fourier}
% \usepackage{newtxmath}
% \usepackage{kpfonts} 

\usepackage[left=2cm, right=2cm, top=2cm, bottom=2cm]{geometry}
\newenvironment{solution}{\begin{proof}[Решение]}{\end{proof}}
\newtheorem*{custom}{}
\newtheorem*{task_1}{Задача 1}
\newtheorem*{1.a}{a}

\title{Домашнее задание - 2}
\author{Падалко Александр, Б05-615}

\begin{document}

\maketitle

\begin{task_1}

	Доказать эффективность TimSort.

\end{task_1}

\begin{solution}
	
	Итак, поддерживаем убывающий стек run-ов.

	Инвариант стека: 

	Для любых трех подряд идущих элементов $Z > Y > X$, $Z > Y + X$.

	Когда инвариант нарушается при добавлении в конец стека очередного
	run-а, мы объединяем с конца $Y$ с минимумом из $X$ и $Z$, пока
	инвариант снова не восстановится.

	\paragraph{a)} Докажем, что это объединение работает за $
	\mathcal{O}(n)$.

	Пока инвариант выполняется у нас поддерживается инвариант, который 
	заставляет расти размеры от верха стека ко дну как минимум как
	числа Фиббоначи.

	Пусть объединяются run-ы с размерами $a_1, a_2, ..., a_i$.
	Тогда, поскольку слияние двух отсортированных run-ов выполняются
	за $\mathcal{O}(len)$ - где, $len$ - суммарная длина объединяемых 
	run-ов. Представим это как то, что $j$-й run объединится 
	$(i - j + 1)$ раз, следовательно от него придет $(i - j + 1) \cdot
	\mathcal{O}(a_j)$.
	Тогда общую асимптотику можем оценить как:

	$\sum_{j = 1}^{i} (i - j + 1)\cdot \mathcal{O}(a_j)$

	Оценим $a_j$. После него, в частности, длины run-ов растут не менее 
	чем как числа Фиббоначи, значит, так как сумма длин всех run-ов 
	не превосходит $n$ и после $a_j$ идет $(i - j + 1)$, то $a_j \cdot 
	F_{i - j + 1} \leq n \implies a_j \leq \frac{n}{F_{i - j + 1}}$.

	Значит $\sum_{j = 1}^{i} (i - j + 1)\cdot \mathcal{O}(a_j) \leq
	\sum_{j = 1}^{i} (i - j + 1) \cdot  \frac{n}{F_{i - j + 1}} = 
	n \cdot \sum_{j = 1}^{i}\frac{i - j + 1}{F_{i - j + 1}} = n\cdot  
	\sum_{j = 1}^{i} \frac{i}{F_i} \leq n\cdot 
	\sum_{j = 1}^{i} \frac{i}{2^i} = n\cdot \mathcal{O}(1) = \mathcal{O}
	(n)$

	\paragraph{b)}

	По сути, в силу того, что при объединении мы объединяем блок
	с большим, то размер объединяемого блока увеличивается хотя бы
	в $2$ раза,
	что говорит о том, что для каждый элемент может участвовать 
	не более чем в $\log_n (n)$ объединениях, а так как для каждого
	конкретного элемента можем считать, что он участвовал только после
	того каждое из этих участий для него закончилось, то суммарно для 
	всех элементов это будет равно $\mathcal{O}(n)$.
	
	

\end{solution}

\end{document}